\documentclass[../../../include/open-logic-section]{subfiles}

\begin{document}

\olfileid{sth}{z}{power}
\olsection{Himpunan Pangkat}

Kita nerusake nganggo aksioma liyane:

\begin{axiom}[Himpunan Pangkat]
Kanggo sembarang himpunan $A$, himpunan
$\Pow{A} = \Setabs{x}{x \subseteq A}$ ana.
\[
	\forall A \exists P \forall x(x \in P \liff (\forall z \in x)z \in A)
\]
\end{axiom}

Panjelasan kang ndhukung iki cukup langsung. Upamane $A$
kawangun ing tataran $S$. Mula kabeh anggotane $A$ wis
kasedhiya sadurunge $S$ (miturut \stagesacc). Dadi kanthi
penalaran kaya ing Pamisahan, saben himpunan bagean $A$
kawangun paling telat ing tataran $S$. Mula kabeh wis
kasedhiya kanggo digabung dadi siji himpunan ing sembarang
tataran sawise $S$. Kita ngerti tataran kaya mangkono ana
amarga $S$ dudu tataran pungkasan (miturut \stagessucc).
Dadi $\Pow{A}$ ana.

Iki akibat kang migunani saka Himpunan Pangkat:

\begin{prop}\label{thm:Products}
Kanggo sembarang himpunan $A, B$, prodhuk Kartesiuse
$A \times B$ ana.
\end{prop}

\begin{proof}
Himpunan $\Pow{\Pow{A \cup B}}$ ana miturut Himpunan Pangkat
lan \olref[sth][z][pairs]{prop:pairsconsequences}. Mula
miturut Pamisahan himpunan iki ana:
\[
	C = \Setabs{z \in \Pow{\Pow{A \cup B}}}{(\exists x \in A)(\exists y \in B) z = \tuple{x, y}}.
\]
Saiki kanggo sembarang $x \in A$ lan $y \in B$,
himpunan $\tuple{x, y}$ ana miturut
\olref[sth][z][pairs]{prop:pairsconsequences}. Kajaba iku,
amarga $x, y \in A \cup B$, kita nduweni
$\{x\}, \{x, y\} \in \Pow{A \cup B}$, lan
$\tuple{x,y} \in \Pow{\Pow{A \cup B}}$. Mula $A \times B = C$.
\end{proof}

Ing bukti iki, Himpunan Pangkat sesambungan karo Pamisahan.
Iki ora nggumunake. Tanpa Pamisahan, Himpunan Pangkat ora
dadi prinsip kang banget \emph{kuwat}. Amarga Pamisahan
ngandharake himpunan bagean endi saka sawijining himpunan
kang ana, lan kanthi mangkono nemtokake sepira ``gedhe''
saben himpunan pangkat.

\begin{prob}
Tuduhna yen kanggo sembarang himpunan $A, B$: (i) himpunan
kabeh relasi kanthi domain $A$ lan range $B$ ana; lan
(ii) himpunan kabeh fungsi saka $A$ menyang $B$ ana.
\end{prob}

\begin{prob}
Upamane $A$ iku himpunan lan $\sim$ relasi ekuivalensi ing $A$.
Buktekna yen himpunan kelas-kelas ekuivalensi miturut $\sim$
ing $A$, yaiku $\equivclass{A}{\sim}$, ana.
\end{prob}

\end{document}
